\documentclass[10pt,a4paper]{amsart}
\usepackage[margin=19mm]{geometry}
\usepackage{amsmath,amssymb,mathtools}
\usepackage[scr=rsfs,cal=euler]{mathalfa}
\usepackage{xfrac}
\usepackage[unicode,hidelinks,bookmarks=false]{hyperref}
\newcommand{\ie}{i.e.\ }
\newcommand{\cf}{cf.\ }
\newcommand{\resp}{respectively\ }
\newcommand{\Subsection}{\subsection}
\providecommand{\nobreakdash}{\nobreak\hbox{-}}
\setlength{\emergencystretch}{2em}
\multlinegap0pt
\usepackage{mathtools}
\usepackage{amssymb}
\usepackage{tikz}
\newtheorem{prop}{Proposition}
\title{Joyce structures from quadratic differentials on the sphere}
\author{Timothy Moy}
\date{Source: arXiv:2509.05275v2 (CC BY 4.0)}
\begin{document}
\maketitle

\noindent\emph{Source and licence.} Joyce structures from quadratic differentials on the sphere. Version arXiv:2509.05275v2 (CC BY 4.0), distributed by arXiv under the Creative Commons Attribution 4.0 International licence, http://creativecommons.org/licenses/by/4.0/, as recorded in the arXiv licence metadata for this exact version and shown on the abstract page https://arxiv.org/abs/2509.05275v2. Full licence notice: \texttt{licenses/arxiv-2509.05275-CC-BY-4.0.txt}.\par\smallskip

\noindent\textit{Editorial scope.} Counted unit: the article's prop prop (source0.tex lines 394--396). The article's complete original proof of it is delivered with it (source0.tex lines 397--441, one \texttt{proof} environment of 45 lines). No mathematics is written, repaired or re-derived: the statement and the proof are the article's own lines, in the article's own notation, and no retained line is substituted or re-flowed.  No further source-local context is needed: every cross-reference of every retained line resolves inside the two delivered spans.  Nothing was omitted from any retained span, so the package carries no omission record.  No retained line contains a citation, so the wrapper carries no bibliography and no bibliography entry.  No retained line contains a figure environment, an image-inclusion command or drawing code, so no image is delivered and the package declares no asset dependency.  Editorial formatting only, in addition: an AMS \texttt{amsart} preamble that restates the article's own declarations and macros used by the retained lines, namely the environments \texttt{prop} on separate counters, together with the standard packages those lines need.  The retained environments print in this document as Proposition 1 in order of appearance, whereas the article prints the same items with the numbering of its own full document.  The complete native source of the article as distributed in the arXiv e-print for this exact version is bundled unchanged as \texttt{sources/184557/source0.tex}.
\begin{prop}
 $\omega$ is a symplectic form.  
\end{prop}

 \begin{proof} 
 It remains to show $\omega$ is non-degenerate.  It is convenient to write $a_{2m_\alpha}^{(\alpha)} \coloneqq w_\alpha$ for each $\alpha = 1,...,N$.  First,  we will show that $\omega$ has a block diagonal decomposition in the coordinate trivialisation
\begin{align}
\bigg\{ \frac{\partial}{\partial a_{i}^{(\alpha)}} \bigg\}_{\alpha = 1,...,N}^{i = 1,...,2m_\alpha} \cup \bigg\{ \frac{\partial}{\partial a_{i}^{(\infty)}} \bigg\}^{i = 0,...,2m_\infty - 7}
\end{align}
for $TM$.  
Then we will show each block is triangular with non-vanishing diagonal entries.  
 At $x = w_\alpha$ we have
\begin{align}\label{order_diff_1}
\frac{\partial y_0}{\partial a_i^{(\alpha)}} = O((x-w_\alpha)^{m_\alpha - \frac{1}{2} - i}),  \quad 
\frac{\partial y_0}{\partial a_j^{(\beta)}} = O((x-w_\alpha)^{m_\alpha - \frac{1}{2}}) 
\end{align}
for $\alpha \ne \beta$ and 
\begin{align}\label{cohomology_coordinate_diff}
\mu_0\bigg(\frac{\partial}{\partial a_i^{(\alpha)}}  \bigg) = \Bigg[ \frac{\partial y_0}{\partial a_i^{(\alpha)}} dx \Bigg] .  
\end{align}
The representative meromorphic $1$-form in (\ref{cohomology_coordinate_diff}) has a pole only at $x = w_\alpha$.  Let $s_\alpha$ be a local coordinate on $\Sigma_0(p)$ at the branch point $x = w_\alpha$ so that $s_\alpha^2 = (x-w_\alpha)$ and $dx = 2s_\alpha ds_\alpha$.  Then 
\begin{align}
\omega \bigg(\frac{\partial }{\partial a_i^{(\alpha)}},  \frac{\partial }{\partial a_j^{(\beta)}} \bigg) = \,  &\underset{x = w_\alpha}{\operatorname{Res}} \bigg( d^{-1} \bigg( \frac{\partial y_0}{\partial a_i^{(\alpha)}}  dx \bigg) \frac{\partial y_0}{\partial a_j^{(\beta)}} dx \bigg) \nonumber +  \,  \underset{x = w_\beta}{\operatorname{Res}} \bigg( d^{-1} \bigg( \frac{\partial y_0}{\partial a_i^{(\alpha)}}  dx \bigg)  \frac{\partial y_0}{\partial a_j^{(\beta)}} dx \bigg) \nonumber \\
= \,  & \underset{s_\alpha = 0}{\operatorname{Res}}(O(s_\alpha^{4m_\alpha+1-2i}) ds_\alpha) +  \underset{s_\beta = 0}{\operatorname{Res}}(O(s_\beta^{4m_\beta+1-2j}) ds_\beta) = 0, 
\end{align}
as $i \le 2m_\alpha$ and  $j \le 2m_\beta$.  Similarly,  at the pole at infinity
\begin{align}\label{order_diff_inf_1}
\frac{\partial y_0}{\partial a_i^{(\alpha)}} = O(x^{-m_\infty+\frac{5}{2}-i}),  \quad \frac{\partial y_0}{\partial a_j^{(\infty)}} = O((x^{j-m_\infty+\frac{5}{2}})).
\end{align}
While,  at a pole $x = w_\alpha$ in $\mathbb{C}$ 
\begin{align}\label{order_diff_inf_2}
\frac{\partial y_0}{\partial a_j^{(\infty)}} = O((x-w_\alpha)^{m_\alpha - \frac{1}{2}}),
\end{align}
so that taking a coordinate $s_\infty^{-2} = x$ about the branch point $x = \infty$ we have 
\begin{align}
\omega\bigg(\frac{\partial }{\partial a_i^{(\alpha)}},  \frac{\partial }{\partial a_j^{(\infty)}} \bigg) = \underset{s_\alpha=0}{\operatorname{Res}}(O(s_\alpha^{4m_\alpha+1-2i})ds_\alpha) + \underset{s_\infty = 0}{\operatorname{Res}}(O(s_\infty^{4m_\infty - 15 - 2j + 2i})ds_\infty)=0,
\end{align}
taking $\alpha = 1,...,N,  i = 1,...,2m_\alpha$ and $j = 0,..., 2m_\infty - 7$.  This gives us a block diagonal decomposition with blocks labelled by $\alpha = 1,...,N$ and $\infty$.  Each block is given by $\omega$ restricted to the vector fields corresponding to the Laurent coefficients at each pole.  For a finite pole $x = w_\alpha$ we have
\begin{align}\label{triangular}
\omega\bigg(\frac{\partial }{\partial a_i^{(\alpha)}},  \frac{\partial }{\partial a_j^{(\alpha)}} \bigg) \nonumber = \,  &\underset{x = w_\alpha}{\operatorname{Res}} \bigg(  d^{-1} \bigg( \frac{\partial y_0}{\partial a_i^{(\alpha)}}  dx \bigg)  \frac{\partial y_0}{\partial a_j^{(\alpha)}} dx \bigg) \\ = \,  &\underset{s_\alpha = 0}{\operatorname{Res}}\bigg( \frac{s_\alpha^{4m_\alpha+1-2i-2j}}{a^{(\alpha)}_{2m_\alpha-1}(2m_\alpha+1-2i)}ds_\alpha+ O(s_\alpha^{4m_\alpha+2-2i-2j})ds_\alpha\bigg),  
\end{align}
which vanishes for $i + j \le 2m_\alpha$.  Since $a_{2m_\alpha-1}^{(\alpha)} \ne 0$ on $M$ (due to the poles of $\phi$ having prescribed orders),  we see that (\ref{triangular}) is non-zero for $i + j = 2m_\alpha + 1$ which shows that the $2m_\alpha \times 2m_\alpha$ matrix with entries (\ref{triangular}) for $i,j = 1,...,2m_\alpha$ is triangular with non-zero diagonal entries (it is triangular in the sense opposite to the usual one). 
Similarly,  for the block corresponding to the pole at infinity,  
\begin{align}\label{triangular_inf}
\omega\bigg(\frac{\partial }{\partial a_i^{(\infty)}},  \frac{\partial }{\partial a_j^{(\infty)}} \bigg) = \,  &\underset{x = \infty}{\operatorname{Res}} \bigg( d^{-1} \bigg( \frac{\partial y_0}{\partial a_i^{(\infty)}}  dx \bigg)  \frac{\partial y_0}{\partial a_j^{(\infty)}} dx \bigg) 
\end{align}
vanishes for $i + j \le 2m_\infty-8$ and is non-zero for $i + j = 2m_\infty - 7$,  so that the last block is triangular.  
Accordingly,  $\omega$ is non-degenerate.  
\end{proof}

\end{document}
